\section{Introduction}
\label{sec:introduction}
Automatic modulation classification (AMC) identifies the modulation format of a
received signal without transmitter-side signaling. It supports cognitive-radio
receivers, spectrum monitoring, and adaptive wireless systems. Convolutional
classifiers applied directly to complex in-phase/quadrature (I/Q) samples and
the RML2016.10a benchmark established a common basis for deep-learning AMC
comparisons \cite{OShea2016Dataset,OShea2016CRR}. Recent surveys describe a
large design space involving signal representations, network structures,
preprocessing, and training procedures \cite{Peng2022Survey,Zhang2022DLAMR,
HuynhThe2021Survey}.

Two design choices recur throughout this literature. The first is \emph{feature
geometry}: a classifier may receive raw I/Q samples or derived amplitude,
phase, polar, constellation, or multi-channel features. The second is
\emph{network capacity}, expressed through width, depth, stages, parameter
count, or operation count. In practice, the two choices are often changed
together. A new representation may be introduced with a new stem, fusion
branch, attention block, or training recipe, while a larger network may also
change the input interface. A reported accuracy difference can therefore
reflect the representation, the capacity used to process it, or the training
and evaluation protocol.

Evidence on the two factors is mixed. Several RML studies report saturation or
limited gains as depth, stages, or other capacity-related choices increase
\cite{West2017Deep,Chen2022SigNet,Harper2023AMC}. For representations, polar
and constellation formulations use transformed signal views
\cite{Peng2019Constellation,Teng2018Polar,Teng2020Polar}. A visualization
study compared I/Q and amplitude/phase inputs under the same LeNet-style and
ResNet-style classifier families over 0--18~dB SNR; the CNN differences were
below one percentage point, whereas an LSTM trained on I/Q failed to learn
\cite{Huang2021Visualizing}. An earlier study also compared amplitude/phase
and I/Q inputs under recurrent and convolutional classifiers
\cite{Rajendran2018Deep}. These results indicate that the effect of input
geometry can depend on the classifier, but they do not determine its size for
the fixed short-window backbone studied here.

Multi-channel and fusion models make attribution more difficult. MCLDNN
combines separate and joint I/Q streams with spatiotemporal processing
\cite{Xu2020MCLDNN}; other methods fuse predictions or learned features across
signal segments or feature views \cite{Zheng2019Fusion,Li2023LightMFFS}.
Recent feature-layer multimodal fusion also combines sequence and image
representations with CNN, LSTM, and attention modules \cite{Wang2024AMFN}.
Such systems are valuable receiver designs, but representation, stem, branch
count, fusion operation, and parameter count may change together. Their results
therefore establish the value of a complete design rather than the isolated
effect of feature geometry.

Lightweight AMC research targets compact models and reduced computational
requirements. MCNet and related efficient CNNs use compact convolutional
structures \cite{HuynhThe2020MCNet,Lin2020HybridNet}; other systems use
multi-feature fusion, embedded-device-oriented time/phase processing, or mobile
Transformer components to obtain accuracy--cost trade-offs
\cite{Cui2023TimePhase,Li2023LightMFFS,Fei2024MobileAmcT,Zheng2023MobileRaT}.
Recent compact examples add cross-channel fusion and attention to an ultra-light
CNN \cite{Wang2024ULNN}, adaptive wavelet extraction to depthwise-separable
convolutions \cite{Quan2024ALWNN}, or residual shrinkage and a CNN--LSTM path
\cite{Suman2025DPDRSN}. These studies report useful parameter, operation, or
runtime comparisons, and several include component or representation ablations,
but the ablated components remain coupled to their proposed architectures.
HFEMCNet is a recent
compact hybrid example that compares FFT-, DWT-, and DCT-based transformations
and reports both component ablations and model-cost comparisons
\cite{Ijaz2026HFEMCNet}. Its ablations vary spectral resolution and pathways,
without providing a frequency-free counterpart, rather than crossing a fixed
feature ladder with a width ladder. These
comparisons are useful for model design, but they commonly change several
architectural ingredients at once. They consequently motivate, rather than
replace, a controlled feature-versus-capacity experiment.

Among the representative works reviewed in Section~\ref{sec:related}, we did
not identify a study that varies feature geometry and network width factorially
on one backbone, under common splits and a common training protocol, with
explicit cost accounting and paired seed- and split-level analysis. The gap is
therefore not the absence of another AMC architecture. It is the absence of a
controlled way to determine whether a change in short-window accuracy is
associated primarily with feature geometry, network capacity, or their
interaction.

This paper addresses that question on short-window RML2016.10a signals. The
corpus, 128-sample observation, backbone family, augmentation, optimizer,
cosine schedule, and evaluation rules are fixed. Four nested feature
configurations, F2, F3, F5, and F7, vary input geometry, while five width
multipliers from $0.50$ to $2.00$ vary the capacity of one depthwise-separable
backbone. Parameters, Core FLOPs, non-MAC FLOPs, and deterministic
preprocessing FLOPs are recorded separately. The ledger supports admissible
cost matching, but the results reported here use nominal width as the capacity
axis and make no accuracy claim from a Core-FLOP-matched pair. A compact second
backbone is reserved for a directional control rather than treated as a second
full benchmark.

H1--H3 are confirmatory hypotheses. H4 and H5 were pre-specified secondary
analyses but are outside the results reported in this manuscript:
\begin{enumerate}
  \item \textbf{H1 (feature effect):} At fixed capacity, moving from F2 to F7 changes final-epoch test accuracy.
  \item \textbf{H2 (capacity effect):} For a fixed feature set, moving from $0.50\times$ to $2.00\times$ width changes accuracy.
  \item \textbf{H3 (interaction):} The accuracy change from added capacity depends on the feature geometry.
  \item \textbf{H4 (SNR regions):} Feature and capacity effects differ between low, middle, and high SNR regions (not reported here).
  \item \textbf{H5 (backbone generality):} The signs of the H1--H3 contrasts persist on a small second-backbone control (not reported here).
\end{enumerate}

The contribution is a controlled experimental design, not a new architecture.
The contributions are as follows:
\begin{enumerate}
  \item A nested feature ladder that makes representation changes explicit
  instead of embedding them in different architectures.
  \item A width ladder on one backbone with one channel-rounding rule, together
  with a cost ledger that separates parameter, Core-FLOP, and preprocessing
  costs.
  \item Predeclared feature, capacity, and interaction contrasts reported with
  seed- and split-level uncertainty; SNR-region and secondary-backbone analyses
  are outside the present results.
  \item Append-only run records, split audits, cost ledgers, and analysis
  configuration files that make the study traceable from its records.
\end{enumerate}

The remainder of this paper is organized as follows. Section~II reviews
related work. Section~III defines the controlled study and the feature and
backbone factors. Section~IV describes the data, splits, training protocol,
and measurements. Section~V reports the planned analyses and results.
Section~VI discusses their implications, Section~VII states the scope and
limitations, and the final sections provide data availability and the
conclusion.

\section{Related Work}
\label{sec:related}
We organize prior work by the factor that each line of research varies and
controls. Table~\ref{tab:literature-matrix} summarizes the representative
works; the categories are analytical groupings rather than claims that every
paper in a category has identical properties.

\begin{table*}[!t]
\caption{Positioning of the proposed study relative to prior work. Selected works; not exhaustive. \emph{Controlled factor attribution} asks whether the focal factor of a category is varied while the backbone, capacity, representation, and training protocol are controlled sufficiently to support attribution. \emph{Independent} means varied by design while the comparison factor is held fixed; \emph{coupled} means changed with other design elements; \emph{mixed} and \emph{partial} indicate heterogeneity among the selected works.}
\label{tab:literature-matrix}
\centering
\footnotesize
\setlength{\tabcolsep}{2pt}
\begin{tabular}{@{}>{\raggedright\arraybackslash}p{0.25\textwidth}>{\raggedright\arraybackslash}p{0.11\textwidth}>{\raggedright\arraybackslash}p{0.13\textwidth}>{\raggedright\arraybackslash}p{0.11\textwidth}>{\raggedright\arraybackslash}p{0.11\textwidth}>{\raggedright\arraybackslash}p{0.13\textwidth}>{\raggedright\arraybackslash}p{0.10\textwidth}@{}}
\toprule
\multicolumn{1}{>{\raggedright\arraybackslash}m{0.25\textwidth}}{Category} &
\multicolumn{1}{>{\raggedright\arraybackslash}m{0.11\textwidth}}{Feature change} &
\multicolumn{1}{>{\raggedright\arraybackslash}m{0.13\textwidth}}{Capacity change} &
\multicolumn{1}{>{\raggedright\arraybackslash}m{0.11\textwidth}}{Same backbone} &
\multicolumn{1}{>{\raggedright\arraybackslash}m{0.11\textwidth}}{Cost reported} &
\multicolumn{1}{>{\raggedright\arraybackslash}m{0.13\textwidth}}{Controlled factor attribution} &
\multicolumn{1}{>{\raggedright\arraybackslash}m{0.10\textwidth}}{Refs.} \\
\midrule
Raw I/Q baselines & Fixed & Mixed & Partial & Partial & Partial & \cite{OShea2016CRR,West2017Deep} \\
Representation and feature fusion & Varied & Mixed & Partial & Partial & Partial & \cite{Peng2019Constellation,Teng2020Polar,Huang2021Visualizing,Wang2024AMFN,Zheng2023MobileRaT,WangGuan2025DiDsNet} \\
Capacity scaling & Mostly fixed & Mixed & Partial & Partial & Partial & \cite{West2017Deep,Chen2022SigNet,Harper2023AMC} \\
Lightweight architectures & Mixed & Varied (coupled) & Partial & Yes & Partial & \cite{HuynhThe2020MCNet,Li2023LightMFFS,Zheng2023MobileRaT,Wang2024ULNN} \\
\textbf{This paper} & \textbf{Varied (independent)} & \textbf{Varied (independent)} & \textbf{Yes} & \textbf{Yes} & \textbf{Yes} & \textbf{--} \\
\bottomrule
\end{tabular}
\end{table*}


\subsection{Raw I/Q baselines}
The RML2016 datasets were generated with GNU Radio and released for machine
learning research on modulation recognition \cite{OShea2016Dataset}. Early
convolutional classifiers operated directly on raw I/Q windows
\cite{OShea2016CRR}. West and O'Shea subsequently varied filters, filter taps,
and depth in related raw-I/Q architectures, finding limited benefit from some
capacity increases \cite{West2017Deep}. These studies establish the raw-I/Q
baseline and motivate capacity analysis, but do not isolate feature geometry.

\subsection{Input representation and feature fusion}
Prior studies use amplitude/phase sequences, constellation diagrams, polar
features, accumulated-polar features, and multi-channel I/Q views
\cite{Rajendran2018Deep,Peng2019Constellation,Teng2018Polar,Teng2020Polar,
Li2023LightMFFS,Xu2020MCLDNN}. The closest controlled comparisons use the same classifier
families for multiple input forms, but they do not combine a feature ladder with
a width ladder and a common paired analysis \cite{Huang2021Visualizing}.
Teng et al. sweep model size for three input forms, but use simulated four-class
data at a single SNR, rather than a paired feature--capacity design
\cite{Teng2020Polar}.
MobileRaT reports I/Q, A/P, and combined-input results at four pruned model
sizes on RML2016.10a, but it provides one value per cell and does not report
cost-matched pairs or paired uncertainty analysis; its size axis is therefore
pruning-based rather than a width ladder \cite{Zheng2023MobileRaT}.
Wang likewise motivates a dual-input design by using I/Q in a low-SNR
path and A/P in a high-SNR path, while examining recurrent depth separately;
the input and capacity factors are not crossed in one factorial experiment
\cite{WangGuan2025DiDsNet}.
Fusion methods additionally combine predictions or features from multiple
segments and views \cite{Zheng2019Fusion,Li2023LightMFFS}. These works show
the breadth of useful signal representations while also illustrating why
feature and architecture effects are difficult to separate.

\subsection{Capacity scaling and saturation}
Studies of depth, width, stages, and model families provide evidence for
saturation on RML data. West and O'Shea examined filter and depth changes
\cite{West2017Deep}; SigNet compared ResNet depths and reported parameter and
training-cost information \cite{Chen2022SigNet}; and Harper et al. related
accuracy to parameterized architectural variants on RML2018.01A with
1024-sample windows \cite{Harper2023AMC}. These results make capacity
saturation a plausible question, but the reported capacity changes are not a
factorial comparison with the same input representation and training protocol.
Width multipliers are a standard way to scale depthwise-separable networks
\cite{Howard2017MobileNets}, while balanced scaling studies motivate treating
width as one explicit design axis rather than as a proxy for every architectural
change \cite{Tan2019EfficientNet}.

\subsection{Lightweight and efficient architectures}
Compact AMC models include efficient CNNs, hybrid recurrent-convolutional
networks, multi-feature fusion networks, and mobile Transformer designs
\cite{HuynhThe2020MCNet,Lin2020HybridNet,Cui2023TimePhase,Li2023LightMFFS,
Fei2024MobileAmcT,Zheng2023MobileRaT,Wang2024ULNN,Quan2024ALWNN,Suman2025DPDRSN}.
They report parameter or operation reductions more often than general AMC
studies, but comparisons across independently designed architectures confound
representation, capacity, and implementation choices. The present study keeps
one primary backbone fixed and changes only the feature channels and width.
Its cost measure is analytical; it does not make a device-level latency or
energy claim.

\subsection{Evaluation protocol, variance, and statistics}
Reported accuracy differences depend on seeds, splits, and evaluation rules.
Benchmark studies show that seed variance can materially change reported
rankings \cite{Henderson2018DeepRL,Bouthillier2021Variance}, and leakage or
undocumented reuse of evaluation information can inflate confidence in
machine-learning conclusions \cite{Kapoor2023Leakage}. Statistical comparisons
of classifiers require care because some resampled-split tests can have
inflated type-I error \cite{Dietterich1998Tests}; Bayesian practical-
equivalence analysis provides a complementary way to describe small
differences \cite{Benavoli2017Time}. Accordingly, this study fixes the split,
seed, checkpoint, cost, and contrast definitions before test access and states
that inference is conditional on the two declared splits.

\subsection{Positioning}
The reviewed works contain valuable representation comparisons, capacity
studies, and efficient architectures. The specific contribution here is their
combination in one controlled design: a nested feature ladder, a width ladder
on a shared backbone, explicit cost accounting, and a predeclared paired
analysis across seeds and two splits. This positioning is deliberately
bounded; it does not claim that the reviewed models are inferior or that the
resulting effects generalize beyond the stated corpus, window, and backbone
scope.
